% Agent 21: Pages 423--426 (PDF pp.\,40--43)
% Sections 5.4 (cont.) and 5.5
% Equations (5.4.1a)--(5.5.9)

%%%%%%%%%%%%%%%%%%%%%%%%%%%%%%%%%%%%%%%%%%%%%%%%%%%%%%%%%%%%%
%%  Section 5.4  Properties of the Kerr Metric Relevant
%%               to Accreting Disks  (continued)
%%%%%%%%%%%%%%%%%%%%%%%%%%%%%%%%%%%%%%%%%%%%%%%%%%%%%%%%%%%%%

Before examining the details of the structure of an accreting disk (\S5.9),
we shall extend our equations of structure from Newtonian theory to general
relativity.  In this extension, we shall assume that the spacetime geometry
outside the hole is that of Kerr, and we shall assume that the disk lies in the
equatorial plane of the Kerr metric.

To treat the structure of such a disk, we shall need a number of properties
of ``direct'' circular orbits in the equatorial plane of the Kerr metric (orbits that
rotate in the same direction as the black hole).  Most of our formulas are taken
from Bardeen, Press and Teukolsky (1972)\cite{Bardeen1972}, or Bardeen (1973)\cite{Bardeen1973} and are readily
derivable from results quoted there.

For simplicity, we shall be splitting formulas into Newtonian limits plus relativistic
corrections, and we shall use the following functions with value unity far
from the hole:
%
\begin{align}
\mathscr{A} &\equiv 1 + a_*^2/r_*^2 + 2a_*^2/r_*^3 , \tag{5.4.1a}\\[4pt]
\mathscr{B} &\equiv 1 + a_*/r_*^{3/2} , \tag{5.4.1b}\\[4pt]
\mathscr{C} &\equiv 1 - 3/r_* + 2a_*/r_*^{3/2} , \tag{5.4.1c}\\[4pt]
\mathscr{D} &\equiv 1 - 2/r_* + a_*^2/r_*^2 , \tag{5.4.1d}\\[4pt]
\mathscr{E} &\equiv 1 + 4a_*^2/r_*^2 - 4a_*^2/r_*^3 + 3a_*^4/r_*^4 , \tag{5.4.1e}\\[4pt]
\mathscr{F} &\equiv 1 - 2a_*/r_*^{3/2} + a_*^2/r_*^2 , \tag{5.4.1f}\\[4pt]
\mathscr{G} &\equiv 1 - 2/r_* + a_*/r_*^{3/2} . \tag{5.4.1g}
\end{align}

Here
\begin{equation}
\mathscr{J} \equiv 1 - 2/r_* + a_* /r_*^{3/2} . \tag{5.4.1h}
\end{equation}
%
\begin{equation}
\varrho \equiv \exp\!\Bigl[-\tfrac{3}{2}\int_{r_{ms}}^{r}\frac{\mathscr{G}\,\mathscr{F}}
{\mathscr{B}\mathscr{C}\,\mathscr{F}^{3/2}}\;\frac{dr_*}{r_{ms}}\Bigr]
= \frac{\mathscr{F}}{\mathscr{G}^{1/2}}\cdot\frac{L_{ms}}{(Mr)^{1/2}}.
\tag{5.4.1i}
\end{equation}
%
\begin{equation}
\varrho = \frac{g^{-3}}{2r_*^{1/2}}\int_{r_{ms}}^{r}
\frac{\mathscr{G}\,\mathscr{F}}
{\mathscr{B}\mathscr{C}\,\mathscr{F}^{3/2}}
\;\frac{dr_*}{r_{ms}}.
\tag{5.4.1j}
\end{equation}

Here $M$ and $a$ are the mass and specific angular momentum of the black hole, $r$ is
radius; $r_*$ and $a_*$ are dimensionless measures of $r$ and $a$
\begin{equation}
r_* = r/M;\qquad a_* = a/M;
\end{equation}

$\tilde{L}_{ms}$ is a constant defined below, and $\tilde{L}$ is a function of $r$ defined below.

The properties of direct circular orbits that we shall need are the following.

\medskip
\noindent
\textbf{(i) Form of Kerr metric in and near equatorial plane}
$\bigl(|\theta - \pi/2| \ll 1\bigr)$:
%
\begin{equation}
ds^2 = -\frac{r^2\Delta}{\mathscr{A}}\,dt^2 + \frac{\mathscr{A}}{r^2}
\bigl(d\phi - \omega\,dt\bigr)^2 + r^2\,d\theta^2 + 2Mra^2\,dz^2 ,
\tag{5.4.2a}
\end{equation}
%
\begin{equation}
\Delta \equiv r^2 - 2Mr + a^2 = r^2\mathscr{D},\qquad
\mathscr{A} = r^4 + r^2a^2 + 2Mra^2 = r^4\,\mathscr{A},
\tag{5.4.2b}
\end{equation}
%
\begin{equation}
\omega = 2Mar/\mathscr{A} = (2Ma/r^3)\,\mathscr{A}^{-1}.
\tag{5.4.2b}
\end{equation}

[We have replaced the usual angular coordinate $\theta$ by
$z = r\cos\theta \simeq r(\theta - \pi/2)$.]
Note that the square root of the determinant of $g_{\alpha\beta}$ is
$(-g)^{1/2} = r$.

\medskip
\noindent
\textbf{(ii) Angular velocity of orbit:}
%
\begin{equation}
\Omega \equiv \frac{d\phi}{dt} = \frac{M^{1/2}}{r^{3/2} + aM^{1/2}}
= \frac{M^{1/2}}{r^{3/2}\,\mathscr{B}}.
\tag{5.4.3}
\end{equation}

\medskip
\noindent
\textbf{(iii) Linear velocity of orbit relative to ``locally nonrotating observer'':}
%
\begin{equation}
V^{(\phi)} = r^{-1}\mathscr{A}^{1/2}\,(\Omega - \omega)
= \frac{M^{1/2}}{r^{1/2}}\;\frac{\mathscr{F}}
{\mathscr{G}^{1/2}\,\mathscr{B}}.
\tag{5.4.4}
\end{equation}

\medskip
\noindent
\textbf{(iv) ``$\gamma$-factor'' corresponding to this linear velocity:}
%
\begin{equation}
\gamma = \bigl(1 - V^{(\phi)2}\bigr)^{-1/2}
= -\frac{\mathscr{B}}{r^{1/2}\mathscr{C}^{1/2}\mathscr{G}}.
\tag{5.4.4b}
\end{equation}

\medskip
\noindent
\textbf{(v) Orthonormal frame attached to an orbiting particle with 4-velocity $u$
(``orbiting frame''):}
%
\begin{align}
\mathbf{e}_{\hat{0}} &= -\frac{\mathscr{B}}{\mathscr{C}^{1/2}}\left(
\frac{r^2}{\mathscr{A}}\right)^{1/2}
\left(\frac{\partial}{\partial t}
+ \Omega\,\frac{\partial}{\partial\phi}\right),
\tag{5.4.5a}\\[6pt]
\mathbf{e}_6 &= r(4/r^2\Delta)^{1/2}
\left(\frac{r^2}{\Delta}\right)^{1/2}
\frac{\partial}{\partial\phi}
+ \gamma V^{(\phi)}\,\mathscr{G}^{1/2}
\frac{\partial}{\partial\phi},\notag\\[3pt]
&= \frac{\mathscr{B}\mathscr{G}^{1/2}}{r^{3/2}\,\mathscr{C}^{1/2}}
\left(\frac{\partial}{\partial t}
+ \omega\,\frac{\partial}{\partial\phi}\right)
+ \frac{\mathscr{G}^{1/2}\,\mathscr{B}}{\mathscr{C}^{1/2}}
\left(\frac{\mathscr{A}}{r}\right)^{1/2}
\frac{\partial}{\partial\phi},
\tag{5.4.5a}\\[6pt]
\mathbf{e}_{\hat{r}} &= \frac{\Delta^{1/2}}{r}\,\frac{\partial}{\partial r}
= \mathscr{D}^{1/2}\,\frac{\partial}{\partial r},
\tag{5.4.5b}\\[6pt]
\mathbf{e}_z &= \frac{\partial}{\partial z}.
\end{align}

\medskip
\noindent
\textbf{(vi) Corresponding orthonormal basis of one-forms:}
%
\begin{align}
\boldsymbol{\omega}^{\hat{0}} &= \frac{\mathscr{G}}{\mathscr{C}^{1/2}}\,dt
- \frac{\mathscr{F}}{\mathscr{C}^{1/2}}\left(\frac{M}{r}\right)^{1/2}
r\,d\phi ,
\tag{5.4.5b}\\[4pt]
\boldsymbol{\omega}^{\hat{3}} &= \frac{\mathscr{B}\mathscr{G}^{1/2}}
{\mathscr{C}^{1/2}}\left(\frac{M}{r}\right)^{1/2}
\frac{1}{r}\,d\phi
- \frac{\mathscr{G}}{\mathscr{C}^{1/2}}\;dt ,
\notag\\[4pt]
\boldsymbol{\omega}^{\hat{3}} &= \mathscr{G}^{-1/2}\,dr_* ,
\tag{5.4.5b}\\[4pt]
\boldsymbol{\omega}^{\hat{3}} &= dz .
\end{align}

\medskip
\noindent
\textbf{(vii) Shear of the congruence of circular, equatorial geodesics} (congruence with
4-velocity $\mathbf{u} = \mathbf{e}_{\hat{0}}$):
%
\begin{equation}
\sigma_{\hat{0}\hat{3}}^{(\text{EG})}
= \sigma_{\hat{3}\hat{0}}^{(\text{EG})} = \frac{3\,M^{1/2}\,\mathscr{G}}
{4\,r^{3/2}\,\mathscr{C}\,\mathscr{B}}\;,
\tag{5.4.6}
\end{equation}
%
all other $\sigma_{\hat{\alpha}\hat{\beta}}^{(\text{EG})}$ vanish.

\medskip
\noindent
\textbf{(viii) Angular momentum per unit mass for circular orbit:}
%
\begin{equation}
\tilde{L} = u_\phi = M^{1/2}\,r^{1/2}\,\mathscr{F}/\mathscr{C}^{1/2}.
\tag{5.4.7a}
\end{equation}

\medskip
\noindent
\textbf{(ix) Energy per unit mass for circular orbit:}
%
\begin{equation}
\tilde{E} = |u_0| = \mathscr{G}/\mathscr{C}^{1/2}.
\tag{5.4.7b}
\end{equation}

\medskip
\noindent
\textbf{(x) Minimum radius for stable circular orbits (``marginally stable'' orbits) $r_{ms}$:}
%
\begin{align}
r_{ms} &- 6Mr_{ms} + 8aM^{1/2}\,r_{ms}^{1/2} - 3a^2 = 0,\notag\\[4pt]
r_{ms} &= M\bigl\{3 + Z_2 - \bigl[(3 - Z_1)(3 + Z_1 + 2Z_2)\bigr]^{1/2}\bigr\},
\tag{5.4.8a}\\[4pt]
Z_1 &= 1 + (1 - a^2/M^2)^{1/3}\bigl[(1 + a/M)^{1/3} + (1 - a/M)^{1/3}\bigr],
\tag{5.4.8b}\\[4pt]
Z_2 &= (3a^2/M^2 + Z_1^2)^{1/2}.
\end{align}

\medskip
\noindent
\textbf{(xi) Angular momentum per unit mass for last stable circular orbit:}
%
\begin{equation}
\tilde{L}_{ms} = M^{1/2}\,r_{ms}^{1/2},\qquad
x = M^{1/2}\,r_{ms}^{1/2}.
\tag{5.4.9}
\end{equation}
%
\begin{equation}
\tilde{L}_{ms} = \frac{2M}{3^{1/2}\,x}\,(3x - 2a).
\end{equation}


%%%%%%%%%%%%%%%%%%%%%%%%%%%%%%%%%%%%%%%%%%%%%%%%%%%%%%%%%%%%%
%%  Section 5.5  Relativistic Model for Disk:
%%               Underlying Assumptions
%%%%%%%%%%%%%%%%%%%%%%%%%%%%%%%%%%%%%%%%%%%%%%%%%%%%%%%%%%%%%

\subsection{Relativistic Model for Disk: Underlying Assumptions}
\label{sec:5.5}

In order to avoid confusion, we shall lay more careful foundations for our
relativistic analysis of disk structure than we did for the Newtonian analysis
of \S5.2.

We shall split the calculation of disk structure into three parts: the radial
structure (\S5.6), the vertical structure (\S\S5.7, 5.8) and the propagation of
radiation from the disk's surface to the observer.

In calculating the disk structure, we shall make the following assumptions
and idealizations:

\begin{enumerate}
\item[(i)] The central plane of the disk coincides with the equatorial plane of the
black hole.

\item[(ii)] The companion star in the binary system has negligible gravitational
influence on the disk structure (assumption valid in inner part of disk;
not at outer edges).

\item[(iii)] The disk is thin; i.e., its proper thickness $2h$, as measured in the orbiting
frame of the disk's matter (same as proper thickness measured by any
other observer who moves in the equatorial plane), satisfies
%
\begin{equation}
h(r) \ll r.
\tag{5.5.1}
\end{equation}

\item[(iv)] The disk is in a quasisteady state.  This assumption can be made precise
in terms of an averaging process.  Pick a particular height $z$ ($|z| \leq h \ll r$)
above the central plane of the disk, and pick a particular quantity $\Psi$
(e.g., $\Psi$ = density or temperature of 4-velocity of gas) to be measured at
height~$z$.  Average $\Psi$ over all $\phi$ (Lie dragging it along $\partial/\partial\phi$ while averaging,
if it is a vector or tensor), average over a proper radial distance of order
$2h$, and average over the time interval required for gas to move inward
a distance $2h$:
%
\begin{equation}
\langle\Psi\rangle = \frac{1}{2\pi\Delta r\,\Delta t}
\int_{r-\epsilon r/2}^{r+\epsilon r/2}
\int_{-\epsilon t/2}^{+\epsilon t/2}
\left(\mathscr{G}^{1/3}\mathscr{B}\right)^{-\epsilon r/2}
\;\Psi\,d\phi\,dt\,dr;
\tag{5.5.2}
\end{equation}
%
$\Delta r = 2h$, $\Delta t = 2h$.  The assumption of a quasisteady state means
that $\langle\Psi\rangle$ is time-independent---i.e.,
%
\begin{equation}
\frac{\partial\langle\Psi\rangle}{\partial t} = 0
\quad\text{if}~\langle\Psi\rangle~\text{is a scalar or tensor},
\tag{5.5.3}
\end{equation}
%
\begin{equation}
\mathscr{L}_{\partial/\partial t}\langle\Psi\rangle = 0
\quad\text{if}~\langle\Psi\rangle~\text{is a vector or tensor},
\end{equation}
%
where $\mathscr{L}_{\partial/\partial t}$ is the Lie derivative along the Killing vector
$\partial/\partial t$.  The
disk structure may be very far from steady state on scales $\lesssim h$; for
example, it may have turbulence, flux reconnection, and flares on such
scales.  Such local violence will not invalidate the analysis that follows.

\item[(v)] When viewed ``macroscopically'' (turbulence removed by the above
averaging process), the gas of the disk moves (very nearly) in direct,
circular, geodesic orbits---i.e., with 4-velocity $\langle\mathbf{u}\rangle \approx \mathbf{e}_{\hat{0}}$.
Superimposed
on this orbital motion is a very small radial flow, produced by viscous
stresses, and an even smaller flow in the vertical direction, as required
by the variation of disk thickness with radius:
%
\begin{equation}
\langle u\rangle = \frac{1}{\bigl[1-(v^{\hat{r}})^2 - (v^{\hat{z}})^2\bigr]^{1/2}}
\bigl[
  \mathbf{e}_{\hat{0}}
+ v^{\hat{r}}\,\mathbf{e}_{\hat{r}}
+ v^{\hat{z}}\,\mathbf{e}_z
\bigr],
\tag{5.5.4}
\end{equation}
%
\begin{equation}
|v^{\hat{r}}(r,z)| \ll V^{(\phi)} \simeq (M/r)^{1/2},
\tag{5.5.5}
\end{equation}
\end{enumerate}

By assuming that the orbital motion of the gas is very nearly geodesic,
we are automatically requiring that the gravitational pull of the hole
dominates over radial pressure gradients and over stresses.  In
order of magnitude the gravitational pull of the hole (``gravitational
acceleration'') is the gradient of the gravitational binding energy,
$(1 - \tilde{E})$; and the acceleration due to pressure gradients and shears is
$\sim (T^{\hat{r}\hat{k}}/\rho_0)_r$.  Thus, we are automatically demanding that the
stress tensor in the fluid's rest frame ($=$ orbiting orthonormal frame)
satisfy
%
\begin{equation}
\frac{T^{\hat{r}\hat{k}}(r,z)}{\rho_0(r,z)} \ll 1 - \tilde{E}(r)
\approx \frac{N}{2}\;\frac{1}{r}\;\frac{M}{r}.
\tag{5.5.6}
\end{equation}

(Here $\approx$ means ``equals in the Newtonian limit''.)

Because the specific internal energy $\Pi$ is approximately equal to
$T^{\hat{k}\hat{k}}/\rho_0$, the assumption of nearly geodesic orbits implies the condition
of ``negligible specific heat'':
%
\begin{equation}
\Pi(r,z) \ll 1 - \tilde{E}(r)
\approx \frac{N}{2}\;\frac{1}{r}\;\frac{M}{r}.
\tag{5.5.7}
\end{equation}

In words---the density of internal energy, $\rho_0\Pi$ (including thermal
energy of gas, energy of turbulent motions, and magnetic-field energy)
is much smaller than the density of gravitational binding energy
$\rho_0(1 - \tilde{E})$.  The mass-averaged radial velocity of the gas

Notice that conditions (5.5.6) and (5.5.7) imply
%
\begin{equation}
\Pi \ll 1, \qquad T_{\hat{r}\hat{k}}/\rho_0 \ll 1.
\tag{5.5.8}
\end{equation}

at all $r$ and $z$---even near the black hole.  This means that, although
general relativistic effects (spacetime curvature; deviations from
Newtonian gravity) are very important near the hole, one can everywhere
ignore special relativistic corrections to the local thermodynamic,
hydrodynamic, and radiative properties of the gas.  Contrast this
situation with the interior of a neutron star and the early stages of the
Universe, where not only is gravity relativistic but so is the matter
$(\Pi \geq 1,\; p/\rho_0 \geq 1,\; T_{\hat{r}\hat{k}}/\rho_0 \geq 1)$.

The mean local rest frame
$[\mathbf{e}_{\hat{0}},\,\mathbf{e}_{\hat{r}},\,\mathbf{e}_z]$;
eqs.\ (5.4.5)].  (iii)  The mean local rest-
frame of the gas
%
\begin{equation}
\bar{\mathbf{e}}_{\hat{0}} = \langle\mathbf{u}\rangle\quad
\text{as given by eq.\ (5.5.4)},
\end{equation}
%
\begin{equation}
\bar{\mathbf{e}}_{\hat{r}} = (\mathbf{e}_{\hat{r}} + v^{\hat{r}}\,\mathbf{e}_{\hat{0}})
/|\mathbf{e}_{\hat{r}} + v^{\hat{r}}\,\mathbf{e}_{\hat{0}}|,
\end{equation}
%
\begin{equation}
\bar{\mathbf{e}}_z = (\mathbf{e}_z + v^{\hat{z}}\,\mathbf{e}_{\hat{0}}
- \hat{v}^{\hat{z}}\,\bar{\mathbf{e}}_{\hat{0}})/
|\mathbf{e}_z + v^{\hat{z}}\,\mathbf{e}_{\hat{0}}
- \hat{v}^{\hat{z}}\,\bar{\mathbf{e}}_{\hat{0}}|.
\tag{5.5.9}
\end{equation}

The mean local rest frame is nearly identical to the orbiting orthonormal frame.
